\documentclass{article}
\usepackage[utf8]{inputenc}
\usepackage[T1]{fontenc}
\usepackage{geometry}
\geometry{
    a4paper,
    left=2.5cm,
    right=2.5cm,
    top=2.5cm,
    bottom=2.5cm
}
\usepackage{graphicx}
\usepackage{float}
\usepackage{hyperref}
\usepackage[english]{babel}
\usepackage{amsmath}
\usepackage{amssymb}
\usepackage{array}
\usepackage{booktabs}
\usepackage{caption}
\usepackage{subcaption}
\usepackage{enumitem}
\usepackage{fancyhdr}
\usepackage{titlesec}
\usepackage{url}
\usepackage{textcomp}
\usepackage{xcolor}
\usepackage{listings}
\usepackage{multirow}
\usepackage{longtable}
\usepackage{tabularx}
\usepackage{times}

% Set up page style
\pagestyle{fancy}
\fancyhf{}
\rfoot{\thepage}
\renewcommand{\headrulewidth}{0pt}

% Title formatting
\titleformat{\chapter}[display]{\normalfont\bfseries}{\chaptertitle\thechapter}{0pt}{\Huge}
\titlespacing*{\chapter}{0pt}{50pt}{40pt}

% Document info
\title{Template Document}
\author{Academic Support Agent}
\date{\today}

\begin{document}
\maketitle

\maketitle
    

    \section{ }
    \begin{frame}{ }
         :
        \begin{itemize}
            \item  
            \item  
        \end{itemize}
    \end{frame}

    \begin{frame}{ }
        \begin{columns}
            \column{.5\textwidth}
             
            
            \column{.5\textwidth}
             
        \end{columns}
    \end{frame}

    \section{ }

    \begin{frame}{ }
        \begin{table}
            \centering
            \begin{tabular}{ccc}
                \hline
                \textbf{ 1} & \textbf{ 2} & \textbf{ 3} \\
                \hline 
                 1 &  2 &  3 \\ 
                 4 &  5 &  6 \\ 
                 7 &  8 &  9 \\ 
                \hline
            \end{tabular}
            \caption{ 1 }
            \label{tab:table1}
        \end{table}
    \end{frame}

    \begin{frame}{ }
        \begin{figure}[H]
            \centering
            \includegraphics[width=.5\textwidth]{example-image-a}
            \caption{ 1 }
            \label{fig:figure1}
        \end{figure}
    \end{frame}

    \begin{frame}{ }
          ( \ref{fig:figure1})\\
          ( \ref{tab:table1})\\
          \cite{vaswani2017attention}\\
        % \footnote{ \textbf{printbibliography} PR Issues  }
          \footfullcite{vaswani2017attention}
    \end{frame}
    
    \lecture{Test lecture one}{lone}
    \begin{frame}{ }
        \begin{theorem}
             1 
        \end{theorem}
        \begin{definition}
             
        \end{definition}
        \begin{corollary}
             
        \end{corollary}
    \end{frame}

    \lecture{Test lecture two}{ltwo}
    \begin{frame}{ }
        \begin{theorem}
             2 
        \end{theorem}
         
        \begin{equation}
            H(x)=-\sum\limits_i^Np(x_i)\log p(x_i).
        \end{equation}
         $H(x)=-\sum\limits_i^Np(x_i)\log p(x_i)$ 
    \end{frame}

    \section{ }

    % Adding the option 'allowframebreaks' allows the contents of the slide to be expanded in more than one slide. 
    % If your references can be listed within one page, may remove the option to rmeove the numbering following the frame title
    \begin{frame}[allowframebreaks]{ }
        % \tiny
        % 
        \printbibliography[heading=none]
    \end{frame}
\end{document}